\chapter{Gesamtinhaltsverzeichnis}

Die folgenden Verzeichnisse geben den Inhalt aller sechsundzwanzig Bände auf Teil- und Kapitelebene wieder. Die Seitenzahlen beziehen sich jeweils auf den einzelnen Band.

\section*{Band 0 \ \textemdash\ \textit{Die Logik des Schalters}}
\addcontentsline{toc}{section}{Band 0 --- Die Logik des Schalters}
\markright{Band 0 --- Die Logik des Schalters}
{\small\color{grau}Grundlagen des Rechnens \ $\cdot$ \ 226 Seiten}
\nopagebreak
{\small
\begin{longtable}{@{}p{0.9cm}p{9.9cm}r@{}}
 & \textit{Vorwort} & 1\\
\multicolumn{3}{@{}l}{\rule{0pt}{3.2ex}\textbf{\color{mittel}I.\ Das Bit}}\\
1 & Was Information ist & 11\\
2 & Zahlen aus Bits & 27\\
3 & Gleitkommazahlen & 41\\
4 & Zeichen, Texte und alles Übrige & 53\\
\multicolumn{3}{@{}l}{\rule{0pt}{3.2ex}\textbf{\color{mittel}II.\ Die Logik}}\\
5 & Die Boolesche Algebra & 67\\
6 & Normalformen und Minimierung & 79\\
7 & Gatter und Schaltnetze & 91\\
\multicolumn{3}{@{}l}{\rule{0pt}{3.2ex}\textbf{\color{mittel}III.\ Das Gedächtnis}}\\
8 & Wie Rückkopplung Gedächtnis erzeugt & 107\\
9 & Endliche Automaten & 117\\
\multicolumn{3}{@{}l}{\rule{0pt}{3.2ex}\textbf{\color{mittel}IV.\ Die Maschine}}\\
10 & Der Von-Neumann-Rechner & 129\\
11 & Maschinensprache und Assembler & 141\\
12 & Der Turm der Abstraktionen & 153\\
\multicolumn{3}{@{}l}{\rule{0pt}{3.2ex}\textbf{\color{mittel}V.\ Das Programm}}\\
13 & Algorithmisches Denken & 167\\
14 & Der Selbstbezug & 179\\
\multicolumn{3}{@{}l}{\rule{0pt}{3.2ex}\textbf{\color{mittel}VI.\ Aktueller Stand}}\\
15 & Wo die Grundlagen heute stehen & 195\\
\multicolumn{3}{@{}l}{\rule{0pt}{3.2ex}\textbf{\color{mittel}VII.\ Anhang}}\\
A & Formelsammlung und Übersicht & 207\\
B & Kommentierte Literatur & 217\\
\end{longtable}}

\section*{Band 1 \ \textemdash\ \textit{Die Ordnung der Daten}}
\addcontentsline{toc}{section}{Band 1 --- Die Ordnung der Daten}
\markright{Band 1 --- Die Ordnung der Daten}
{\small\color{grau}Algorithmen und Datenstrukturen \ $\cdot$ \ 216 Seiten}
\nopagebreak
{\small
\begin{longtable}{@{}p{0.9cm}p{9.9cm}r@{}}
 & \textit{Vorwort} & 1\\
\multicolumn{3}{@{}l}{\rule{0pt}{3.2ex}\textbf{\color{mittel}I.\ Grundlagen}}\\
1 & Aufwand messen & 11\\
2 & Felder, Listen und ihre Verwandten & 25\\
\multicolumn{3}{@{}l}{\rule{0pt}{3.2ex}\textbf{\color{mittel}II.\ Ordnen und Suchen}}\\
3 & Die elementaren Sortierverfahren & 37\\
4 & Die schnellen Sortierverfahren & 47\\
5 & Suchen und Auswählen & 63\\
\multicolumn{3}{@{}l}{\rule{0pt}{3.2ex}\textbf{\color{mittel}III.\ Bäume}}\\
6 & Binäre Suchbäume & 75\\
7 & Balancierte Bäume & 85\\
8 & Heaps und Prioritätswarteschlangen & 95\\
\multicolumn{3}{@{}l}{\rule{0pt}{3.2ex}\textbf{\color{mittel}IV.\ Streuspeicherung}}\\
9 & Hashing & 107\\
\multicolumn{3}{@{}l}{\rule{0pt}{3.2ex}\textbf{\color{mittel}V.\ Graphen}}\\
10 & Graphen darstellen und durchmustern & 123\\
11 & Kürzeste Wege & 137\\
12 & Spannbäume und Flüsse & 149\\
\multicolumn{3}{@{}l}{\rule{0pt}{3.2ex}\textbf{\color{mittel}VI.\ Entwurfsmuster}}\\
13 & Teile und herrsche, und die gierigen Verfahren & 163\\
14 & Dynamische Programmierung & 175\\
\multicolumn{3}{@{}l}{\rule{0pt}{3.2ex}\textbf{\color{mittel}VII.\ Aktueller Stand}}\\
15 & Wo die Algorithmik heute steht & 187\\
A & Übersicht und Formelsammlung & 197\\
B & Kommentierte Literatur & 207\\
\end{longtable}}

\section*{Band 2 \ \textemdash\ \textit{Die Grenze der Maschine}}
\addcontentsline{toc}{section}{Band 2 --- Die Grenze der Maschine}
\markright{Band 2 --- Die Grenze der Maschine}
{\small\color{grau}Berechenbarkeit und formale Sprachen \ $\cdot$ \ 176 Seiten}
\nopagebreak
{\small
\begin{longtable}{@{}p{0.9cm}p{9.9cm}r@{}}
 & \textit{Vorwort} & 1\\
\multicolumn{3}{@{}l}{\rule{0pt}{3.2ex}\textbf{\color{mittel}I.\ Sprachen und Automaten}}\\
1 & Wörter, Sprachen und das Abzählen & 11\\
2 & Endliche Automaten & 19\\
3 & Reguläre Ausdrücke und ihre Grenze & 27\\
4 & Kontextfreie Sprachen & 37\\
\multicolumn{3}{@{}l}{\rule{0pt}{3.2ex}\textbf{\color{mittel}II.\ Die Maschine}}\\
5 & Die Turingmaschine & 49\\
6 & Die universelle Maschine & 59\\
7 & Die Church-Turing-These & 67\\
\multicolumn{3}{@{}l}{\rule{0pt}{3.2ex}\textbf{\color{mittel}III.\ Die Grenze}}\\
8 & Das Halteproblem & 79\\
9 & Reduktionen und der Satz von Rice & 87\\
10 & Aufzählbarkeit & 97\\
11 & Grade der Unlösbarkeit & 105\\
\multicolumn{3}{@{}l}{\rule{0pt}{3.2ex}\textbf{\color{mittel}IV.\ Jenseits der Informatik}}\\
12 & Unentscheidbares in der Mathematik & 117\\
13 & Gödels Unvollständigkeitssätze & 129\\
14 & Kolmogorow-Komplexität & 139\\
\multicolumn{3}{@{}l}{\rule{0pt}{3.2ex}\textbf{\color{mittel}V.\ Aktueller Stand}}\\
15 & Wo die Theorie heute steht & 151\\
\multicolumn{3}{@{}l}{\rule{0pt}{3.2ex}\textbf{\color{mittel}VI.\ Anhang}}\\
A & Übersicht und Formelsammlung & 161\\
B & Kommentierte Literatur & 169\\
\end{longtable}}

\section*{Band 3 \ \textemdash\ \textit{Der Preis der Berechnung}}
\addcontentsline{toc}{section}{Band 3 --- Der Preis der Berechnung}
\markright{Band 3 --- Der Preis der Berechnung}
{\small\color{grau}Komplexitätstheorie \ $\cdot$ \ 166 Seiten}
\nopagebreak
{\small
\begin{longtable}{@{}p{0.9cm}p{9.9cm}r@{}}
 & \textit{Vorwort} & 1\\
\multicolumn{3}{@{}l}{\rule{0pt}{3.2ex}\textbf{\color{mittel}I.\ Ressourcen messen}}\\
1 & Zeit und Platz & 11\\
2 & Die Klasse P & 19\\
\multicolumn{3}{@{}l}{\rule{0pt}{3.2ex}\textbf{\color{mittel}II.\ Nichtdeterminismus}}\\
3 & Die Klasse NP & 31\\
4 & Reduktionen und NP-Vollständigkeit & 39\\
5 & Der Zoo der vollständigen Probleme & 49\\
\multicolumn{3}{@{}l}{\rule{0pt}{3.2ex}\textbf{\color{mittel}III.\ Struktur}}\\
6 & coNP und die Hierarchie & 61\\
7 & Platz als Ressource & 69\\
8 & Rechnen fast ohne Speicher & 77\\
9 & Hierarchiesätze und die Mitte & 83\\
\multicolumn{3}{@{}l}{\rule{0pt}{3.2ex}\textbf{\color{mittel}IV.\ Wege um die Härte}}\\
10 & Näherungsverfahren & 93\\
11 & Zufall als Ressource & 103\\
12 & Parametrisierte Komplexität & 113\\
13 & Härte als Werkzeug & 123\\
\multicolumn{3}{@{}l}{\rule{0pt}{3.2ex}\textbf{\color{mittel}V.\ Grenzen der Methode}}\\
14 & Warum die Frage offen ist & 133\\
\multicolumn{3}{@{}l}{\rule{0pt}{3.2ex}\textbf{\color{mittel}VI.\ Aktueller Stand}}\\
15 & Wo die Theorie heute steht & 145\\
\multicolumn{3}{@{}l}{\rule{0pt}{3.2ex}\textbf{\color{mittel}VII.\ Anhang}}\\
A & Übersicht und Formelsammlung & 153\\
B & Kommentierte Literatur & 159\\
\end{longtable}}

\section*{Band 4 \ \textemdash\ \textit{Die Anatomie der Geschwindigkeit}}
\addcontentsline{toc}{section}{Band 4 --- Die Anatomie der Geschwindigkeit}
\markright{Band 4 --- Die Anatomie der Geschwindigkeit}
{\small\color{grau}Rechnerarchitektur \ $\cdot$ \ 166 Seiten}
\nopagebreak
{\small
\begin{longtable}{@{}p{0.9cm}p{9.9cm}r@{}}
 & \textit{Vorwort} & 1\\
\multicolumn{3}{@{}l}{\rule{0pt}{3.2ex}\textbf{\color{mittel}I.\ Der Datenpfad}}\\
1 & Wie ein Prozessor aufgebaut ist & 11\\
2 & Der Befehlssatz als Vertrag & 19\\
\multicolumn{3}{@{}l}{\rule{0pt}{3.2ex}\textbf{\color{mittel}II.\ Das Fließband}}\\
3 & Fließbandverarbeitung & 31\\
4 & Konflikte im Fließband & 39\\
5 & Sprungvorhersage & 47\\
6 & Mehrere Befehle gleichzeitig & 55\\
\multicolumn{3}{@{}l}{\rule{0pt}{3.2ex}\textbf{\color{mittel}III.\ Das Gedächtnis}}\\
7 & Die Speicherhierarchie & 67\\
8 & Cachebewusst programmieren & 77\\
9 & Virtueller Speicher & 87\\
\multicolumn{3}{@{}l}{\rule{0pt}{3.2ex}\textbf{\color{mittel}IV.\ Vielfalt und Parallelität}}\\
10 & Vektorbefehle & 97\\
11 & Mehrere Kerne & 105\\
12 & Grafikprozessoren und Spezialbausteine & 113\\
\multicolumn{3}{@{}l}{\rule{0pt}{3.2ex}\textbf{\color{mittel}V.\ Wo es hakt}}\\
13 & Energie als Grenze & 123\\
14 & Wenn die Abstraktion undicht wird & 131\\
\multicolumn{3}{@{}l}{\rule{0pt}{3.2ex}\textbf{\color{mittel}VI.\ Aktueller Stand}}\\
15 & Wo die Architektur heute steht & 143\\
\multicolumn{3}{@{}l}{\rule{0pt}{3.2ex}\textbf{\color{mittel}VII.\ Anhang}}\\
A & Übersicht und Zahlen & 151\\
B & Kommentierte Literatur & 159\\
\end{longtable}}

\section*{Band 5 \ \textemdash\ \textit{Der Wettlauf der Prozesse}}
\addcontentsline{toc}{section}{Band 5 --- Der Wettlauf der Prozesse}
\markright{Band 5 --- Der Wettlauf der Prozesse}
{\small\color{grau}Betriebssysteme und Nebenläufigkeit \ $\cdot$ \ 186 Seiten}
\nopagebreak
{\small
\begin{longtable}{@{}p{0.9cm}p{9.9cm}r@{}}
 & \textit{Vorwort} & 1\\
\multicolumn{3}{@{}l}{\rule{0pt}{3.2ex}\textbf{\color{mittel}I.\ Der Kern}}\\
1 & Was ein Betriebssystem ist & 11\\
2 & Unterbrechungen & 23\\
\multicolumn{3}{@{}l}{\rule{0pt}{3.2ex}\textbf{\color{mittel}II.\ Prozesse}}\\
3 & Prozesse und Fäden & 35\\
4 & Ablaufplanung & 47\\
5 & Speicherverwaltung & 57\\
\multicolumn{3}{@{}l}{\rule{0pt}{3.2ex}\textbf{\color{mittel}III.\ Nebenläufigkeit}}\\
6 & Wettläufe und wechselseitiger Ausschluss & 67\\
7 & Semaphore und Monitore & 79\\
8 & Verklemmungen & 91\\
9 & Die klassischen Aufgaben & 99\\
\multicolumn{3}{@{}l}{\rule{0pt}{3.2ex}\textbf{\color{mittel}IV.\ Ein- und Ausgabe}}\\
10 & Geräte und Treiber & 111\\
11 & Dateisysteme & 121\\
12 & Wenn der Strom ausfällt & 131\\
\multicolumn{3}{@{}l}{\rule{0pt}{3.2ex}\textbf{\color{mittel}V.\ Größere Zusammenhänge}}\\
13 & Virtualisierung und Container & 143\\
14 & Schutz und Isolation & 153\\
15 & Der Stand im Jahr 2026 & 163\\
\multicolumn{3}{@{}l}{\rule{0pt}{3.2ex}\textbf{\color{mittel}VI.\ Anhang}}\\
A & Zahlen und Kenngrößen & 171\\
B & Landkarte des Bandes & 173\\
C & Glossar & 175\\
D & Zeittafel & 177\\
E & Kommentierte Quellen & 179\\
1 & 1 384 721          6                1 204 & 556\\
2 & 1 102 993          7                2 000 & 0\\
3 & 1 551 208          8                1 317 & 842\\
4 & 1 299 441          9                1 468 & 15\\
5 & 1 733 760          10               1 090 & 377\\
\end{longtable}}

\section*{Band 6 \ \textemdash\ \textit{Die Verwandlung des Textes}}
\addcontentsline{toc}{section}{Band 6 --- Die Verwandlung des Textes}
\markright{Band 6 --- Die Verwandlung des Textes}
{\small\color{grau}Programmiersprachen und Übersetzerbau \ $\cdot$ \ 188 Seiten}
\nopagebreak
{\small
\begin{longtable}{@{}p{0.9cm}p{9.9cm}r@{}}
 & \textit{Vorwort} & 1\\
\multicolumn{3}{@{}l}{\rule{0pt}{3.2ex}\textbf{\color{mittel}I.\ Vom Text zur Struktur}}\\
1 & Was ein Übersetzer tut & 11\\
2 & Lexikalische Analyse & 21\\
3 & Syntaktische Analyse & 31\\
4 & Namen und Bedeutung & 41\\
\multicolumn{3}{@{}l}{\rule{0pt}{3.2ex}\textbf{\color{mittel}II.\ Die Sprachen selbst}}\\
5 & Typsysteme & 53\\
6 & Speicher und Lebensdauer & 63\\
7 & Abstraktionsmittel & 73\\
\multicolumn{3}{@{}l}{\rule{0pt}{3.2ex}\textbf{\color{mittel}III.\ Vom Baum zum Befehl}}\\
8 & Zwischendarstellungen & 85\\
9 & Optimierung & 95\\
10 & Codeerzeugung & 107\\
11 & Binden und Laden & 117\\
\multicolumn{3}{@{}l}{\rule{0pt}{3.2ex}\textbf{\color{mittel}IV.\ Andere Wege}}\\
12 & Interpreter und virtuelle Maschinen & 127\\
13 & Übersetzen zur Laufzeit & 139\\
14 & Übersetzer in der Praxis & 149\\
15 & Der Stand im Jahr 2026 & 161\\
\multicolumn{3}{@{}l}{\rule{0pt}{3.2ex}\textbf{\color{mittel}V.\ Anhang}}\\
A & Zahlen und Kenngrößen & 171\\
B & Landkarte des Bandes & 175\\
C & Glossar & 177\\
D & Zeittafel & 179\\
E & Kommentierte Quellen & 181\\
10 & und 10         100 & 20\\
\end{longtable}}

\section*{Band 7 \ \textemdash\ \textit{Das Maß der Nachricht}}
\addcontentsline{toc}{section}{Band 7 --- Das Maß der Nachricht}
\markright{Band 7 --- Das Maß der Nachricht}
{\small\color{grau}Informationstheorie \ $\cdot$ \ 158 Seiten}
\nopagebreak
{\small
\begin{longtable}{@{}p{0.9cm}p{9.9cm}r@{}}
 & \textit{Vorwort} & 1\\
\multicolumn{3}{@{}l}{\rule{0pt}{3.2ex}\textbf{\color{mittel}I.\ Das Maß}}\\
1 & Was Information ist & 9\\
2 & Entropie & 17\\
3 & Was zwei Größen voneinander wissen & 27\\
\multicolumn{3}{@{}l}{\rule{0pt}{3.2ex}\textbf{\color{mittel}II.\ Verdichten}}\\
4 & Codes & 35\\
5 & Der Quellencodierungssatz & 43\\
6 & Verfahren zur Verdichtung & 51\\
7 & Wörterbücher und Modelle & 61\\
8 & Wenn man etwas verlieren darf & 69\\
\multicolumn{3}{@{}l}{\rule{0pt}{3.2ex}\textbf{\color{mittel}III.\ Der Kanal}}\\
9 & Kanäle und Kapazität & 81\\
10 & Der Kanalcodierungssatz & 89\\
11 & Fehler erkennen und beheben & 95\\
12 & Codes, die die Grenze erreichen & 105\\
\multicolumn{3}{@{}l}{\rule{0pt}{3.2ex}\textbf{\color{mittel}IV.\ Weiterungen}}\\
13 & Die Information einer einzelnen Folge & 117\\
14 & Dieselbe Größe in anderen Fächern & 125\\
15 & Der Stand im Jahr 2026 & 133\\
\multicolumn{3}{@{}l}{\rule{0pt}{3.2ex}\textbf{\color{mittel}V.\ Anhang}}\\
A & Formeln und Zahlen & 141\\
B & Landkarte des Bandes & 145\\
C & Glossar & 147\\
D & Zeittafel & 149\\
E & Kommentierte Quellen & 151\\
\end{longtable}}

\section*{Band 8 \ \textemdash\ \textit{Die Mathematik des Misstrauens}}
\addcontentsline{toc}{section}{Band 8 --- Die Mathematik des Misstrauens}
\markright{Band 8 --- Die Mathematik des Misstrauens}
{\small\color{grau}Kryptographie \ $\cdot$ \ 174 Seiten}
\nopagebreak
{\small
\begin{longtable}{@{}p{0.9cm}p{9.9cm}r@{}}
 & \textit{Vorwort} & 1\\
\multicolumn{3}{@{}l}{\rule{0pt}{3.2ex}\textbf{\color{mittel}I.\ Grundlagen}}\\
1 & Was geschützt werden soll & 11\\
2 & Klassische Verfahren und ihr Bruch & 19\\
3 & Der Übergang zum Aufwand & 29\\
\multicolumn{3}{@{}l}{\rule{0pt}{3.2ex}\textbf{\color{mittel}II.\ Symmetrische Verfahren}}\\
4 & Blockchiffren & 41\\
5 & Betriebsarten & 51\\
6 & Stromchiffren und Zufall & 61\\
7 & Hashfunktionen & 69\\
8 & Nachrichten authentifizieren & 79\\
\multicolumn{3}{@{}l}{\rule{0pt}{3.2ex}\textbf{\color{mittel}III.\ Asymmetrische Verfahren}}\\
9 & Rechnen mit Resten & 89\\
10 & RSA & 97\\
11 & Diffie-Hellman & 107\\
12 & Elliptische Kurven & 117\\
13 & Signaturen und Zertifikate & 125\\
\multicolumn{3}{@{}l}{\rule{0pt}{3.2ex}\textbf{\color{mittel}IV.\ Praxis und Zukunft}}\\
14 & Wo es tatsächlich schiefgeht & 137\\
15 & Der Stand im Jahr 2026 & 147\\
\multicolumn{3}{@{}l}{\rule{0pt}{3.2ex}\textbf{\color{mittel}V.\ Anhang}}\\
A & Zahlen und Empfehlungen & 157\\
B & Landkarte des Bandes & 161\\
C & Glossar & 163\\
D & Zeittafel & 165\\
E & Kommentierte Quellen & 167\\
\end{longtable}}

\section*{Band 9 \ \textemdash\ \textit{Die Einigung der Getrennten}}
\addcontentsline{toc}{section}{Band 9 --- Die Einigung der Getrennten}
\markright{Band 9 --- Die Einigung der Getrennten}
{\small\color{grau}Netzwerke und verteilte Systeme \ $\cdot$ \ 166 Seiten}
\nopagebreak
{\small
\begin{longtable}{@{}p{0.9cm}p{9.9cm}r@{}}
 & \textit{Vorwort} & 1\\
\multicolumn{3}{@{}l}{\rule{0pt}{3.2ex}\textbf{\color{mittel}I.\ Das Netz}}\\
1 & Warum Schichten & 11\\
2 & Die unteren Schichten & 21\\
3 & Adressen und Wegewahl & 31\\
4 & Transport & 41\\
5 & Namen und Anwendungen & 51\\
\multicolumn{3}{@{}l}{\rule{0pt}{3.2ex}\textbf{\color{mittel}II.\ Verteiltes Rechnen}}\\
6 & Was es schwer macht & 63\\
7 & Zeit und Ordnung & 71\\
8 & Kopien und Konsistenz & 79\\
9 & Konsens & 87\\
10 & Wenn Knoten lügen & 95\\
\multicolumn{3}{@{}l}{\rule{0pt}{3.2ex}\textbf{\color{mittel}III.\ Bauen im Großen}}\\
11 & Daten verteilen & 105\\
12 & Nachrichten und Ströme & 113\\
13 & Ausfälle und Betrieb & 123\\
14 & Leistung im Verteilten & 131\\
15 & Der Stand im Jahr 2026 & 141\\
\multicolumn{3}{@{}l}{\rule{0pt}{3.2ex}\textbf{\color{mittel}IV.\ Anhang}}\\
A & Zahlen und Formeln & 149\\
B & Landkarte des Bandes & 153\\
C & Glossar & 155\\
D & Zeittafel & 157\\
E & Kommentierte Quellen & 159\\
\end{longtable}}

\section*{Band 10 \ \textemdash\ \textit{Das Gedächtnis der Systeme}}
\addcontentsline{toc}{section}{Band 10 --- Das Gedächtnis der Systeme}
\markright{Band 10 --- Das Gedächtnis der Systeme}
{\small\color{grau}Datenbanken und Information Retrieval \ $\cdot$ \ 162 Seiten}
\nopagebreak
{\small
\begin{longtable}{@{}p{0.9cm}p{9.9cm}r@{}}
 & \textit{Vorwort} & 1\\
\multicolumn{3}{@{}l}{\rule{0pt}{3.2ex}\textbf{\color{mittel}I.\ Modell und Sprache}}\\
1 & Wozu ein Datenbanksystem & 11\\
2 & Das relationale Modell & 21\\
3 & SQL & 29\\
4 & Entwurf und Normalformen & 37\\
\multicolumn{3}{@{}l}{\rule{0pt}{3.2ex}\textbf{\color{mittel}II.\ Wie es innen arbeitet}}\\
5 & Speicherung & 47\\
6 & Indexstrukturen & 55\\
7 & Anfrageausführung & 65\\
8 & Anfrageoptimierung & 73\\
\multicolumn{3}{@{}l}{\rule{0pt}{3.2ex}\textbf{\color{mittel}III.\ Gleichzeitigkeit und Dauerhaftigkeit}}\\
9 & Transaktionen und Isolation & 85\\
10 & Nebenläufigkeitskontrolle & 93\\
11 & Wiederanlauf & 103\\
\multicolumn{3}{@{}l}{\rule{0pt}{3.2ex}\textbf{\color{mittel}IV.\ Finden statt Abfragen}}\\
12 & Der invertierte Index & 113\\
13 & Rangfolge & 121\\
14 & Verbindungen und Vektoren & 129\\
15 & Der Stand im Jahr 2026 & 139\\
\multicolumn{3}{@{}l}{\rule{0pt}{3.2ex}\textbf{\color{mittel}V.\ Anhang}}\\
A & Zahlen und Formeln & 147\\
B & Landkarte des Bandes & 149\\
C & Glossar & 151\\
D & Zeittafel & 153\\
E & Kommentierte Quellen & 155\\
\end{longtable}}

\section*{Band 11 \ \textemdash\ \textit{Das Muster der Verbindungen}}
\addcontentsline{toc}{section}{Band 11 --- Das Muster der Verbindungen}
\markright{Band 11 --- Das Muster der Verbindungen}
{\small\color{grau}Graphen und Netzwerkanalyse \ $\cdot$ \ 164 Seiten}
\nopagebreak
{\small
\begin{longtable}{@{}p{0.9cm}p{9.9cm}r@{}}
 & \textit{Vorwort} & 1\\
\multicolumn{3}{@{}l}{\rule{0pt}{3.2ex}\textbf{\color{mittel}I.\ Grundlagen}}\\
1 & Begriffe und Darstellungen & 11\\
2 & Durchmusterung und Zusammenhang & 21\\
3 & Wege und Flüsse & 29\\
\multicolumn{3}{@{}l}{\rule{0pt}{3.2ex}\textbf{\color{mittel}II.\ Struktur messen}}\\
4 & Zentralität & 39\\
5 & Lokale Struktur & 47\\
6 & Gemeinschaften & 55\\
7 & Hierarchie und Ordnung & 63\\
\multicolumn{3}{@{}l}{\rule{0pt}{3.2ex}\textbf{\color{mittel}III.\ Modelle}}\\
8 & Zufallsgraphen & 73\\
9 & Kleine Welten und schwere Enden & 81\\
10 & Spektrale Graphentheorie & 91\\
\multicolumn{3}{@{}l}{\rule{0pt}{3.2ex}\textbf{\color{mittel}IV.\ Anwendung}}\\
11 & Prozesse auf Netzen & 101\\
12 & Das Gehirn als Netzwerk & 109\\
13 & Große Graphen berechnen & 119\\
14 & Fallstricke & 129\\
15 & Der Stand im Jahr 2026 & 139\\
\multicolumn{3}{@{}l}{\rule{0pt}{3.2ex}\textbf{\color{mittel}V.\ Anhang}}\\
A & Formeln und Zahlen & 147\\
B & Landkarte des Bandes & 151\\
C & Glossar & 153\\
D & Zeittafel & 155\\
E & Kommentierte Quellen & 157\\
\end{longtable}}

\section*{Band 12 \ \textemdash\ \textit{Der Raum der Möglichkeiten}}
\addcontentsline{toc}{section}{Band 12 --- Der Raum der Möglichkeiten}
\markright{Band 12 --- Der Raum der Möglichkeiten}
{\small\color{grau}Symbolische künstliche Intelligenz \ $\cdot$ \ 164 Seiten}
\nopagebreak
{\small
\begin{longtable}{@{}p{0.9cm}p{9.9cm}r@{}}
 & \textit{Vorwort} & 1\\
\multicolumn{3}{@{}l}{\rule{0pt}{3.2ex}\textbf{\color{mittel}I.\ Suchen}}\\
1 & Zeichen und Zustandsräume & 11\\
2 & Suche ohne Vorwissen & 19\\
3 & Heuristische Suche & 27\\
4 & Suche gegen einen Gegner & 37\\
5 & Lokale Suche & 47\\
\multicolumn{3}{@{}l}{\rule{0pt}{3.2ex}\textbf{\color{mittel}II.\ Nebenbedingungen und Logik}}\\
6 & Nebenbedingungen & 57\\
7 & Aussagenlogik und SAT & 65\\
8 & Prädikatenlogik & 73\\
9 & Logisches Programmieren & 81\\
\multicolumn{3}{@{}l}{\rule{0pt}{3.2ex}\textbf{\color{mittel}III.\ Wissen und Handeln}}\\
10 & Wissen darstellen & 93\\
11 & Planen & 101\\
12 & Unsicherheit & 109\\
13 & Nichtmonotones Schließen & 117\\
\multicolumn{3}{@{}l}{\rule{0pt}{3.2ex}\textbf{\color{mittel}IV.\ Bilanz}}\\
14 & Was schiefging und was blieb & 127\\
15 & Der Stand im Jahr 2026 & 137\\
\multicolumn{3}{@{}l}{\rule{0pt}{3.2ex}\textbf{\color{mittel}V.\ Anhang}}\\
A & Formeln und Zahlen & 147\\
B & Landkarte des Bandes & 151\\
C & Glossar & 153\\
D & Zeittafel & 155\\
E & Kommentierte Quellen & 157\\
\end{longtable}}

\section*{Band 13 \ \textemdash\ \textit{Der Schluss aus dem Beispiel}}
\addcontentsline{toc}{section}{Band 13 --- Der Schluss aus dem Beispiel}
\markright{Band 13 --- Der Schluss aus dem Beispiel}
{\small\color{grau}Maschinelles Lernen \ $\cdot$ \ 170 Seiten}
\nopagebreak
{\small
\begin{longtable}{@{}p{0.9cm}p{9.9cm}r@{}}
 & \textit{Vorwort} & 1\\
\multicolumn{3}{@{}l}{\rule{0pt}{3.2ex}\textbf{\color{mittel}I.\ Grundlagen}}\\
1 & Was Lernen heißt & 11\\
2 & Anpassen und Verallgemeinern & 21\\
3 & Lineare Modelle & 29\\
4 & Bäume und Ensembles & 37\\
5 & Abstände und Kerne & 45\\
6 & Wahrscheinlichkeitsbasierte Verfahren & 55\\
7 & Unüberwachtes Lernen & 65\\
\multicolumn{3}{@{}l}{\rule{0pt}{3.2ex}\textbf{\color{mittel}II.\ Neuronale Netze}}\\
8 & Netze und Rückpropagierung & 79\\
9 & Optimierung im Großen & 87\\
10 & Architekturen & 95\\
11 & Sprachmodelle & 103\\
\multicolumn{3}{@{}l}{\rule{0pt}{3.2ex}\textbf{\color{mittel}III.\ Praxis und Grenzen}}\\
12 & Bewertung & 113\\
13 & Fallstricke & 123\\
14 & Erklärbarkeit und Kausalität & 133\\
15 & Der Stand im Jahr 2026 & 143\\
\multicolumn{3}{@{}l}{\rule{0pt}{3.2ex}\textbf{\color{mittel}IV.\ Anhang}}\\
A & Formeln und Zahlen & 153\\
B & Landkarte des Bandes & 157\\
C & Glossar & 159\\
D & Zeittafel & 161\\
E & Kommentierte Quellen & 163\\
4 & Matrizen, & 9\\
7 & Verdichten = & 11\\
9 & verteiltes & 9\\
11 & Einbettungen & 7\\
12 & Suche als & 9\\
\end{longtable}}

\section*{Band 14 \ \textemdash\ \textit{Das Rechnen des Gehirns}}
\addcontentsline{toc}{section}{Band 14 --- Das Rechnen des Gehirns}
\markright{Band 14 --- Das Rechnen des Gehirns}
{\small\color{grau}Computational Neuroscience \ $\cdot$ \ 174 Seiten}
\nopagebreak
{\small
\begin{longtable}{@{}p{0.9cm}p{9.9cm}r@{}}
 & \textit{Vorwort} & 1\\
\multicolumn{3}{@{}l}{\rule{0pt}{3.2ex}\textbf{\color{mittel}I.\ Die Einzelzelle}}\\
1 & Membran und Potential & 11\\
2 & Das Aktionspotential & 19\\
3 & Vereinfachte Modelle & 27\\
\multicolumn{3}{@{}l}{\rule{0pt}{3.2ex}\textbf{\color{mittel}II.\ Synapsen und Plastizität}}\\
4 & Synaptische Übertragung & 37\\
5 & Lernen an der Synapse & 47\\
6 & Dendriten als Rechenelemente & 57\\
\multicolumn{3}{@{}l}{\rule{0pt}{3.2ex}\textbf{\color{mittel}III.\ Kodierung}}\\
7 & Der neuronale Code & 69\\
8 & Information im Nervensystem & 79\\
9 & Rezeptive Felder & 89\\
\multicolumn{3}{@{}l}{\rule{0pt}{3.2ex}\textbf{\color{mittel}IV.\ Netzwerke}}\\
10 & Netzwerkdynamik & 99\\
11 & Die Architektur des Kortex & 109\\
12 & Das Konnektom & 117\\
\multicolumn{3}{@{}l}{\rule{0pt}{3.2ex}\textbf{\color{mittel}V.\ Anwendung und Grenzen}}\\
13 & Computational Psychiatry & 131\\
14 & Was Modelle erklären & 141\\
15 & Der Stand im Jahr 2026 & 149\\
\multicolumn{3}{@{}l}{\rule{0pt}{3.2ex}\textbf{\color{mittel}VI.\ Anhang}}\\
A & Formeln und Zahlen & 157\\
B & Landkarte des Bandes & 161\\
C & Glossar & 163\\
D & Zeittafel & 165\\
E & Kommentierte Quellen & 167\\
5 & Nebenläufigkeit & 10\\
7 & Information, & 8\\
9 & verteilte & 12\\
11 & Graphen, & 12\\
12 & Bayes-Netze, & 13\\
13 & Lernen, Vorher- & 5\\
\end{longtable}}

\section*{Band 15 \ \textemdash\ \textit{Die Auslöschung der Irrwege}}
\addcontentsline{toc}{section}{Band 15 --- Die Auslöschung der Irrwege}
\markright{Band 15 --- Die Auslöschung der Irrwege}
{\small\color{grau}Quanteninformatik \ $\cdot$ \ 156 Seiten}
\nopagebreak
{\small
\begin{longtable}{@{}p{0.9cm}p{9.9cm}r@{}}
 & \textit{Vorwort} & 1\\
\multicolumn{3}{@{}l}{\rule{0pt}{3.2ex}\textbf{\color{mittel}I.\ Grundlagen}}\\
1 & Das Qubit & 11\\
2 & Gatter und Schaltungen & 19\\
3 & Verschränkung & 27\\
\multicolumn{3}{@{}l}{\rule{0pt}{3.2ex}\textbf{\color{mittel}II.\ Die Verfahren}}\\
4 & Interferenz als Rechenprinzip & 37\\
5 & Die Quanten-Fouriertransformation & 43\\
6 & Shors Algorithmus & 51\\
7 & Grovers Algorithmus & 59\\
\multicolumn{3}{@{}l}{\rule{0pt}{3.2ex}\textbf{\color{mittel}III.\ Grenzen und Anschluss}}\\
8 & Was nicht geht & 69\\
9 & Quanteninformationstheorie & 77\\
10 & Quantenkryptographie & 85\\
\multicolumn{3}{@{}l}{\rule{0pt}{3.2ex}\textbf{\color{mittel}IV.\ Umsetzung}}\\
11 & Dekohärenz und Rauschen & 95\\
12 & Fehlerkorrektur & 103\\
13 & Bauformen & 111\\
\multicolumn{3}{@{}l}{\rule{0pt}{3.2ex}\textbf{\color{mittel}V.\ Bilanz}}\\
14 & Was gilt und was behauptet wird & 121\\
15 & Der Stand im Jahr 2026 & 129\\
\multicolumn{3}{@{}l}{\rule{0pt}{3.2ex}\textbf{\color{mittel}VI.\ Anhang}}\\
A & Formeln und Zahlen & 137\\
B & Landkarte des Bandes & 143\\
C & Glossar & 145\\
D & Zeittafel & 147\\
E & Kommentierte Quellen & 149\\
2 & Berechenbarkeit & 8\\
3 & Komplexitäts- & 8\\
4 & Gatter, & 2\\
7 & Entropie, & 9\\
8 & Kryptographie & 10\\
13 & Optimierung, & 15\\
\end{longtable}}

\section*{Begleitband \ \textemdash\ \textit{Der seltsame Kreis}}
\addcontentsline{toc}{section}{Begleitband --- Der seltsame Kreis}
\markright{Begleitband --- Der seltsame Kreis}
{\small\color{grau}Selbstbezug, Universalität und Grenze \ $\cdot$ \ 124 Seiten}
\nopagebreak
{\small
\begin{longtable}{@{}p{0.9cm}p{9.9cm}r@{}}
 & \textit{Vorwort} & 1\\
\multicolumn{3}{@{}l}{\rule{0pt}{3.2ex}\textbf{\color{mittel}I.\ Die Bewegung}}\\
1 & Die Diagonale & 11\\
2 & Der Selbstbezug & 17\\
3 & Universalität & 25\\
\multicolumn{3}{@{}l}{\rule{0pt}{3.2ex}\textbf{\color{mittel}II.\ Die Grenzen}}\\
4 & Gödel & 35\\
5 & Turing & 41\\
6 & Wo die Diagonale versagt & 47\\
\multicolumn{3}{@{}l}{\rule{0pt}{3.2ex}\textbf{\color{mittel}III.\ Der Kreis in der Technik}}\\
7 & Der Übersetzer, der sich selbst übersetzt & 57\\
8 & Das Programm, das sich selbst kopiert & 65\\
9 & Das System, das sich selbst beschreibt & 73\\
\multicolumn{3}{@{}l}{\rule{0pt}{3.2ex}\textbf{\color{mittel}IV.\ Der Kreis im Denken}}\\
10 & Das Gehirn, das sich selbst modelliert & 83\\
11 & Die Grenze anderer Bauart & 91\\
12 & Was die Reihe gezeigt hat & 97\\
\multicolumn{3}{@{}l}{\rule{0pt}{3.2ex}\textbf{\color{mittel}V.\ Anhang}}\\
A & Die Bewegung im Überblick & 107\\
B & Landkarte & 109\\
C & Glossar & 113\\
D & Zeittafel & 115\\
E & Kommentierte Quellen & 117\\
\end{longtable}}

\section*{Band 16 \ \textemdash\ \textit{Die Maschine, die rechnen lernte}}
\addcontentsline{toc}{section}{Band 16 --- Die Maschine, die rechnen lernte}
\markright{Band 16 --- Die Maschine, die rechnen lernte}
{\small\color{grau}Eine Geschichte des maschinellen Rechnens \ $\cdot$ \ 146 Seiten}
\nopagebreak
{\small
\begin{longtable}{@{}p{0.9cm}p{9.9cm}r@{}}
 & \textit{Vorwort} & 1\\
\multicolumn{3}{@{}l}{\rule{0pt}{3.2ex}\textbf{\color{mittel}I.\ Vor der Maschine}}\\
1 & Zählen als Kulturtechnik & 11\\
2 & Die mechanischen Rechenmaschinen & 17\\
3 & Die Idee des Programms & 23\\
\multicolumn{3}{@{}l}{\rule{0pt}{3.2ex}\textbf{\color{mittel}II.\ Die Maschine entsteht}}\\
4 & Rechnen als Verwaltungstechnik & 33\\
5 & Turing und die Grundlagenkrise & 41\\
6 & Der Krieg als Beschleuniger & 47\\
7 & Das gespeicherte Programm & 55\\
\multicolumn{3}{@{}l}{\rule{0pt}{3.2ex}\textbf{\color{mittel}III.\ Die Maschine wird zum Betrieb}}\\
8 & Großrechner und Stapelverarbeitung & 63\\
9 & Transistor, Chip und ein Gesetz, das keines ist & 69\\
10 & Mittlere Datentechnik & 77\\
\multicolumn{3}{@{}l}{\rule{0pt}{3.2ex}\textbf{\color{mittel}IV.\ Die Maschine wird persönlich}}\\
11 & Der Mikroprozessor & 89\\
12 & Der Heimcomputer & 97\\
13 & Arbeitsplatzrechner und Unix & 105\\
14 & Die Vernetzung & 113\\
\multicolumn{3}{@{}l}{\rule{0pt}{3.2ex}\textbf{\color{mittel}V.\ Bilanz}}\\
15 & Der Stand im Jahr 2026 & 123\\
\multicolumn{3}{@{}l}{\rule{0pt}{3.2ex}\textbf{\color{mittel}VI.\ Anhang}}\\
A & Zahlen und Kenngrößen & 131\\
B & Landkarte des Bandes & 135\\
C & Glossar & 137\\
D & Kommentierte Quellen & 139\\
\end{longtable}}

\section*{Band 17 \ \textemdash\ \textit{Das Werk vieler Hände}}
\addcontentsline{toc}{section}{Band 17 --- Das Werk vieler Hände}
\markright{Band 17 --- Das Werk vieler Hände}
{\small\color{grau}Softwaretechnik \ $\cdot$ \ 150 Seiten}
\nopagebreak
{\small
\begin{longtable}{@{}p{0.9cm}p{9.9cm}r@{}}
 & \textit{Vorwort} & 1\\
\multicolumn{3}{@{}l}{\rule{0pt}{3.2ex}\textbf{\color{mittel}I.\ Der Gegenstand}}\\
1 & Was Software anders macht & 11\\
2 & Größe und ihre Folgen & 17\\
3 & Der Lebenszyklus und seine Kosten & 25\\
\multicolumn{3}{@{}l}{\rule{0pt}{3.2ex}\textbf{\color{mittel}II.\ Entwurf}}\\
4 & Module & 35\\
5 & Architektur & 43\\
6 & Schnittstellen und Verträge & 51\\
\multicolumn{3}{@{}l}{\rule{0pt}{3.2ex}\textbf{\color{mittel}III.\ Sicherstellung}}\\
7 & Testen & 61\\
8 & Statische Prüfung & 69\\
9 & Formale Verfahren & 77\\
\multicolumn{3}{@{}l}{\rule{0pt}{3.2ex}\textbf{\color{mittel}IV.\ Der Betrieb der Entwicklung}}\\
10 & Versionsverwaltung und Integration & 87\\
11 & Vorgehensmodelle & 95\\
12 & Schätzen und Messen & 103\\
\multicolumn{3}{@{}l}{\rule{0pt}{3.2ex}\textbf{\color{mittel}V.\ Die Zeit}}\\
13 & Wartung und technische Schuld & 113\\
14 & Ablösung und Migration & 119\\
15 & Der Stand im Jahr 2026 & 127\\
\multicolumn{3}{@{}l}{\rule{0pt}{3.2ex}\textbf{\color{mittel}VI.\ Anhang}}\\
A & Formeln und Kenngrößen & 135\\
B & Landkarte und Glossar & 139\\
C & Zeittafel & 141\\
D & Kommentierte Quellen & 143\\
\end{longtable}}

\section*{Band 18 \ \textemdash\ \textit{Die Übersetzung der Absicht}}
\addcontentsline{toc}{section}{Band 18 --- Die Übersetzung der Absicht}
\markright{Band 18 --- Die Übersetzung der Absicht}
{\small\color{grau}Mensch-Maschine-Interaktion \ $\cdot$ \ 142 Seiten}
\nopagebreak
{\small
\begin{longtable}{@{}p{0.9cm}p{9.9cm}r@{}}
 & \textit{Vorwort} & 1\\
\multicolumn{3}{@{}l}{\rule{0pt}{3.2ex}\textbf{\color{mittel}I.\ Der Mensch}}\\
1 & Wahrnehmung und Aufmerksamkeit & 11\\
2 & Gedächtnis und Belastung & 19\\
3 & Die Gesetze der Bewegung und der Auswahl & 25\\
\multicolumn{3}{@{}l}{\rule{0pt}{3.2ex}\textbf{\color{mittel}II.\ Die Kluft}}\\
4 & Die beiden Übersetzungen & 35\\
5 & Mentale Modelle & 43\\
6 & Fehler & 49\\
\multicolumn{3}{@{}l}{\rule{0pt}{3.2ex}\textbf{\color{mittel}III.\ Der Entwurf}}\\
7 & Sichtbarkeit und Angebot & 59\\
8 & Die Zeit & 67\\
9 & Rückgängigmachen und Vertrauen & 73\\
\multicolumn{3}{@{}l}{\rule{0pt}{3.2ex}\textbf{\color{mittel}IV.\ Die Prüfung}}\\
10 & Bedienzeiten vorhersagen & 81\\
11 & Empirische Prüfung & 89\\
12 & Barrierefreiheit & 95\\
\multicolumn{3}{@{}l}{\rule{0pt}{3.2ex}\textbf{\color{mittel}V.\ Die Praxis}}\\
13 & Muster und Gegenmuster & 105\\
14 & Ist es schlechter geworden? & 111\\
15 & Der Stand im Jahr 2026 & 119\\
\multicolumn{3}{@{}l}{\rule{0pt}{3.2ex}\textbf{\color{mittel}VI.\ Anhang}}\\
A & Formeln und Kenngrößen & 127\\
B & Landkarte und Glossar & 131\\
C & Zeittafel & 133\\
D & Kommentierte Quellen & 135\\
\end{longtable}}

\section*{Band 19 \ \textemdash\ \textit{Der schwächste Punkt}}
\addcontentsline{toc}{section}{Band 19 --- Der schwächste Punkt}
\markright{Band 19 --- Der schwächste Punkt}
{\small\color{grau}Sicherheit \ $\cdot$ \ 150 Seiten}
\nopagebreak
{\small
\begin{longtable}{@{}p{0.9cm}p{9.9cm}r@{}}
 & \textit{Vorwort} & 1\\
\multicolumn{3}{@{}l}{\rule{0pt}{3.2ex}\textbf{\color{mittel}I.\ Die Denkweise}}\\
1 & Die Asymmetrie & 11\\
2 & Das Angreifermodell & 19\\
3 & Angriffsfläche und Vertrauensgrenzen & 27\\
\multicolumn{3}{@{}l}{\rule{0pt}{3.2ex}\textbf{\color{mittel}II.\ Wo Systeme brechen}}\\
4 & Speicherfehler & 37\\
5 & Eingaben und Einschleusung & 45\\
6 & Rechte und Identität & 51\\
\multicolumn{3}{@{}l}{\rule{0pt}{3.2ex}\textbf{\color{mittel}III.\ Das Umfeld}}\\
7 & Konfiguration und Voreinstellungen & 61\\
8 & Die Lieferkette & 67\\
9 & Der Mensch & 75\\
\multicolumn{3}{@{}l}{\rule{0pt}{3.2ex}\textbf{\color{mittel}IV.\ Was sich verrät}}\\
10 & Seitenkanäle & 85\\
11 & Kryptographie in der Anwendung & 91\\
12 & Protokolle und Zustände & 99\\
\multicolumn{3}{@{}l}{\rule{0pt}{3.2ex}\textbf{\color{mittel}V.\ Verteidigung}}\\
13 & Härten & 109\\
14 & Erkennen und Reagieren & 117\\
\multicolumn{3}{@{}l}{\rule{0pt}{3.2ex}\textbf{\color{mittel}VI.\ Bilanz}}\\
15 & Der Stand im Jahr 2026 & 127\\
\multicolumn{3}{@{}l}{\rule{0pt}{3.2ex}\textbf{\color{mittel}VII.\ Anhang}}\\
A & Formeln und Kenngrößen & 135\\
B & Landkarte und Glossar & 139\\
C & Zeittafel & 141\\
D & Kommentierte Quellen & 143\\
\end{longtable}}

\section*{Band 20 \ \textemdash\ \textit{Der zweite Entwurf}}
\addcontentsline{toc}{section}{Band 20 --- Der zweite Entwurf}
\markright{Band 20 --- Der zweite Entwurf}
{\small\color{grau}Programmieren in Python \ $\cdot$ \ 139 Seiten}
\nopagebreak
{\small
\begin{longtable}{@{}p{0.9cm}p{9.9cm}r@{}}
 & \textit{Vorwort} & 1\\
\multicolumn{3}{@{}l}{\rule{0pt}{3.2ex}\textbf{\color{mittel}I.\ Bevor man schreibt}}\\
1 & Was die Aufgabe wirklich ist & 9\\
2 & Daten zuerst & 15\\
3 & Der Ablauf & 25\\
\multicolumn{3}{@{}l}{\rule{0pt}{3.2ex}\textbf{\color{mittel}II.\ Der erste Entwurf}}\\
4 & Der naheliegende Weg & 35\\
5 & Zerlegung & 43\\
6 & Zustand & 51\\
\multicolumn{3}{@{}l}{\rule{0pt}{3.2ex}\textbf{\color{mittel}III.\ Das Handwerk}}\\
7 & Fehlersuche als Verfahren & 61\\
8 & Fremden Code lesen & 69\\
9 & Prüfen & 75\\
\multicolumn{3}{@{}l}{\rule{0pt}{3.2ex}\textbf{\color{mittel}IV.\ Python im Besonderen}}\\
10 & Was Python anders macht & 83\\
11 & Die Werkzeuge der Sprache & 89\\
12 & Wo Python weh tut & 97\\
\multicolumn{3}{@{}l}{\rule{0pt}{3.2ex}\textbf{\color{mittel}V.\ Der zweite Entwurf}}\\
13 & Umschreiben & 105\\
14 & Wann aufhören & 111\\
15 & Der Stand im Jahr 2026 & 119\\
\multicolumn{3}{@{}l}{\rule{0pt}{3.2ex}\textbf{\color{mittel}VI.\ Anhang}}\\
A & Regeln und Kenngrößen & 127\\
B & Landkarte und Glossar & 131\\
C & Kommentierte Quellen & 133\\
\end{longtable}}

\section*{Band 21 \ \textemdash\ \textit{Was darunter liegt}}
\addcontentsline{toc}{section}{Band 21 --- Was darunter liegt}
\markright{Band 21 --- Was darunter liegt}
{\small\color{grau}Programmieren in C \ $\cdot$ \ 132 Seiten}
\nopagebreak
{\small
\begin{longtable}{@{}p{0.9cm}p{9.9cm}r@{}}
 & \textit{Vorwort} & 1\\
\multicolumn{3}{@{}l}{\rule{0pt}{3.2ex}\textbf{\color{mittel}I.\ Die Maschine wird sichtbar}}\\
1 & Was verborgen war & 9\\
2 & Speicher & 15\\
3 & Zeiger und Besitz & 23\\
\multicolumn{3}{@{}l}{\rule{0pt}{3.2ex}\textbf{\color{mittel}II.\ Der erste Entwurf in C}}\\
4 & Zeichenketten & 33\\
5 & Fehler ohne Ausnahmen & 39\\
6 & Datenstrukturen von Hand & 45\\
\multicolumn{3}{@{}l}{\rule{0pt}{3.2ex}\textbf{\color{mittel}III.\ Die Fallen}}\\
7 & Undefiniertes Verhalten & 55\\
8 & Ganzzahlen und Gleitkomma & 63\\
9 & Der Übersetzer & 69\\
\multicolumn{3}{@{}l}{\rule{0pt}{3.2ex}\textbf{\color{mittel}IV.\ Das Handwerk}}\\
10 & Werkzeuge & 77\\
11 & Module und Schnittstellen & 83\\
12 & Fehlersuche in C & 89\\
\multicolumn{3}{@{}l}{\rule{0pt}{3.2ex}\textbf{\color{mittel}V.\ Das Urteil}}\\
13 & Der zweite Entwurf in C & 99\\
14 & Wann C und wann nicht & 105\\
15 & Der Stand im Jahr 2026 & 111\\
\multicolumn{3}{@{}l}{\rule{0pt}{3.2ex}\textbf{\color{mittel}VI.\ Anhang}}\\
A & Regeln und Kenngrößen & 119\\
B & Landkarte und Glossar & 123\\
C & Kommentierte Quellen & 125\\
\end{longtable}}

\section*{Band 22 \ \textemdash\ \textit{Die Kunst des Weglassens}}
\addcontentsline{toc}{section}{Band 22 --- Die Kunst des Weglassens}
\markright{Band 22 --- Die Kunst des Weglassens}
{\small\color{grau}Programmieren in Go \ $\cdot$ \ 118 Seiten}
\nopagebreak
{\small
\begin{longtable}{@{}p{0.9cm}p{9.9cm}r@{}}
 & \textit{Vorwort} & 1\\
\multicolumn{3}{@{}l}{\rule{0pt}{3.2ex}\textbf{\color{mittel}I.\ Die Auslassungen}}\\
1 & Was fehlt und warum & 9\\
2 & Fehler als Werte & 15\\
3 & Schnittstellen & 21\\
\multicolumn{3}{@{}l}{\rule{0pt}{3.2ex}\textbf{\color{mittel}II.\ Nebenläufigkeit}}\\
4 & Goroutinen & 29\\
5 & Kanäle & 35\\
6 & Wo es schiefgeht & 41\\
\multicolumn{3}{@{}l}{\rule{0pt}{3.2ex}\textbf{\color{mittel}III.\ Der Entwurf}}\\
7 & Pakete & 49\\
8 & Zusammensetzung & 55\\
9 & Der Nullwert & 61\\
\multicolumn{3}{@{}l}{\rule{0pt}{3.2ex}\textbf{\color{mittel}IV.\ Das Handwerk}}\\
10 & Werkzeuge & 69\\
11 & Prüfen & 75\\
12 & Leistung und Speicherbereinigung & 81\\
\multicolumn{3}{@{}l}{\rule{0pt}{3.2ex}\textbf{\color{mittel}V.\ Das Urteil}}\\
13 & Der zweite Entwurf in Go & 89\\
14 & Wann Go und wann nicht & 95\\
15 & Der Stand im Jahr 2026 & 99\\
\multicolumn{3}{@{}l}{\rule{0pt}{3.2ex}\textbf{\color{mittel}VI.\ Anhang}}\\
A & Regeln und Kenngrößen & 107\\
B & Landkarte und Glossar & 109\\
C & Kommentierte Quellen & 111\\
\end{longtable}}

\section*{Band 23 \ \textemdash\ \textit{Vom Programm zum Werk}}
\addcontentsline{toc}{section}{Band 23 --- Vom Programm zum Werk}
\markright{Band 23 --- Vom Programm zum Werk}
{\small\color{grau}Programme, die zusammenwirken \ $\cdot$ \ 122 Seiten}
\nopagebreak
{\small
\begin{longtable}{@{}p{0.9cm}p{9.9cm}r@{}}
 & \textit{Vorwort} & 1\\
\multicolumn{3}{@{}l}{\rule{0pt}{3.2ex}\textbf{\color{mittel}I.\ Das Werk als Ganzes}}\\
1 & Programm und Werk & 9\\
2 & Die Naht & 15\\
3 & Verträge & 21\\
\multicolumn{3}{@{}l}{\rule{0pt}{3.2ex}\textbf{\color{mittel}II.\ Über die Grenze}}\\
4 & Zeit & 29\\
5 & Text & 35\\
6 & Zahlen & 41\\
\multicolumn{3}{@{}l}{\rule{0pt}{3.2ex}\textbf{\color{mittel}III.\ Leistung im Großen}}\\
7 & Messen im Verbund & 49\\
8 & Wo die Zeit hingeht & 55\\
9 & Zwischenspeicher & 61\\
\multicolumn{3}{@{}l}{\rule{0pt}{3.2ex}\textbf{\color{mittel}IV.\ Der Betrieb}}\\
10 & Fehler im Verbund & 69\\
11 & Beobachtbarkeit & 75\\
12 & Ausfall und Wiederanlauf & 81\\
\multicolumn{3}{@{}l}{\rule{0pt}{3.2ex}\textbf{\color{mittel}V.\ Das Urteil}}\\
13 & Der zweite Entwurf im Großen & 89\\
14 & Wann teilen und wann nicht & 95\\
15 & Der Stand im Jahr 2026 & 101\\
\multicolumn{3}{@{}l}{\rule{0pt}{3.2ex}\textbf{\color{mittel}VI.\ Anhang}}\\
A & Regeln und Kenngrößen & 109\\
B & Landkarte und Glossar & 113\\
C & Kommentierte Quellen & 115\\
\end{longtable}}

\section*{Band 24 \ \textemdash\ \textit{Die Kunst}}
\addcontentsline{toc}{section}{Band 24 --- Die Kunst}
\markright{Band 24 --- Die Kunst}
{\small\color{grau}Zwanzig Fälle und ihre Begründung \ $\cdot$ \ 108 Seiten}
\nopagebreak
{\small
\begin{longtable}{@{}p{0.9cm}p{9.9cm}r@{}}
 & \textit{Vorwort} & 1\\
\multicolumn{3}{@{}l}{\rule{0pt}{3.2ex}\textbf{\color{mittel}I.\ Die Daten}}\\
1 & Der Wert ohne Einheit & 7\\
2 & Die Funktion mit acht Argumenten & 11\\
3 & Die verfrühte Abstraktion & 15\\
4 & Das Feld, das fehlen darf & 19\\
\multicolumn{3}{@{}l}{\rule{0pt}{3.2ex}\textbf{\color{mittel}II.\ Der Ablauf}}\\
5 & Die Schleife, die zu viel fragt & 25\\
6 & Die Prüfung, die verschwindet & 29\\
7 & Der Puffer ohne Grenze & 33\\
8 & Die Goroutine, die nie endet & 37\\
9 & Der Kanal, der zu viel puffert & 41\\
\multicolumn{3}{@{}l}{\rule{0pt}{3.2ex}\textbf{\color{mittel}III.\ Die Struktur}}\\
10 & Die Hierarchie, die nicht mehr passt & 47\\
11 & Der Zwischenspeicher ohne Ablaufzeit & 51\\
12 & Die Zeit, die falsch gewählt war & 55\\
13 & Die Wiederholung ohne Kennung & 59\\
14 & Die Meldung ohne Zusammenhang & 63\\
\multicolumn{3}{@{}l}{\rule{0pt}{3.2ex}\textbf{\color{mittel}IV.\ Wann man es lässt}}\\
15 & Der Code, der schon gut war & 69\\
16 & Die Optimierung, die bleibt & 73\\
17 & Der doppelte Code, der bleibt & 77\\
18 & Die doppelte Prüfung & 81\\
19 & Der Dienst, der bleibt & 85\\
20 & Das Werkzeug, das bleibt & 89\\
\multicolumn{3}{@{}l}{\rule{0pt}{3.2ex}\textbf{\color{mittel}V.\ Anhang}}\\
A & Die zwanzig Fälle & 95\\
B & Landkarte und Glossar & 99\\
C & Kommentierte Quellen & 101\\
\multicolumn{3}{@{}l}{\rule{0pt}{3.2ex}\textbf{\color{mittel}I.\ 1–4            Die Daten        20,}}\\
\multicolumn{3}{@{}l}{\rule{0pt}{3.2ex}\textbf{\color{mittel}II.\ 5–9            Der Ablauf       21,}}\\
\multicolumn{3}{@{}l}{\rule{0pt}{3.2ex}\textbf{\color{mittel}III.\ 10–14          Die Struktur     17,}}\\
C & 5           Band & 21\\
\end{longtable}}

